\chapter{Guessing It: Primal--Dual Active Sets and Block Pivoting}
\label{ch:guessing}

The dual method of Chapter~\ref{ch:walking} reaches the active set one constraint
at a time. That is its defining feature and, at moderate $n$, its principal cost: a
budget-plus-bounds problem in a few hundred variables takes on the order of a
hundred outer iterations, each a handful of matrix--vector products.

The alternative is to guess the whole active set at once, solve the working-set
subproblem it induces, and repair the guess from the signs that come back. This
chapter develops that idea, proves it terminates finitely in the bound-constrained
case, and shows precisely why the proof does not survive the general constraint
class. The last point is the chapter's main contribution to the reader's judgement:
it is a case study in \emph{recognising that a theorem's hypothesis is structural},
not a technicality to be patched.

\section{The iteration}

Given a candidate working set $\Active$, Proposition~\ref{prop:sub} gives the
solution directly:
\begin{equation}
  \bigl(C_\Active\T G^{-1} C_\Active\bigr)\lambda_\Active
   = b_\Active - C_\Active\T x_u,
  \qquad
  x = x_u + G^{-1}C_\Active \lambda_\Active,
  \qquad x_u = G^{-1}a,
  \label{eq:pdassolve}
\end{equation}
with $\lambda_i = 0$ off $\Active$. Both solves go through one Cholesky
factorisation of $G$, computed once for the whole run, and one of the $k\times k$
dual Hessian, computed per repair.

The repair rule reads the two sign conditions
\eqref{eq:dfeas}--\eqref{eq:pfeas} as instructions:
\index{primal--dual active set}
\begin{equation}
  \Active' = \{1,\dots,m_{\mathrm{eq}}\} \;\cup\;
  \bigl\{\, i > m_{\mathrm{eq}} \;:\;
    (i \in \Active \text{ and } \lambda_i \ge 0)
    \ \text{ or }\
    (c_i\T x < b_i) \,\bigr\}.
  \label{eq:repair}
\end{equation}
In words: an active constraint whose multiplier has gone negative does not belong
and leaves; an inactive constraint that is violated at $x$ does belong and joins;
equalities are always in. The starting guess is the equalities together with
whatever the unconstrained minimiser $x_u$ violates---which is already the answer
when it violates nothing. If $\Active' = \Active$ the KKT conditions hold to the
tolerance used and~\eqref{eq:pdassolve} supplies stationarity by construction.

The whole set is exchanged at once, and in practice this converges in a handful of
repairs \emph{almost independently of $n$}. It trades arithmetic, which is abundant
at these sizes, for iterations, which are not---a trade Chapter~\ref{ch:measurement}
argues is the right one below a few hundred variables and the wrong one above a few
thousand.

\begin{example}[Why the guess is usually good]
For a long-only minimum-variance problem, $x_u = G^{-1}a$ is the unconstrained
minimum-variance portfolio, which typically has a few dozen negative weights out of
several hundred assets. Those are exactly the bounds the first guess activates,
and most of them are correct. The repairs correct the rest.
\end{example}

\section{Two readings of the same iteration}

The iteration above has two lineages, and each contributes something.

\paragraph{Semismooth Newton.} Hintermüller, Ito and Kunisch \cite{hintermuller2002}
show that the primal--dual active set strategy is a semismooth Newton method applied
to the complementarity condition written as a nonsmooth equation
$\lambda - \max(0, \lambda - \sigma(c_i\T x - b_i)) = 0$. This explains the observed
behaviour: Newton's method converges in few iterations and its iteration count does
not scale with the dimension. It also explains the failure mode---Newton's method
without globalisation may not converge at all.

\paragraph{Block principal pivoting.} Specialise to bound constraints, so that by
Chapter~\ref{ch:geometry} the problem is $\mathrm{LCP}(A,-b)$. Toggling one index
between free and bound is a \emph{principal pivot}; toggling several at once is a
\emph{block} principal pivot \cite{judice1994,kim2011}. The repair
rule~\eqref{eq:repair} is exactly a block pivot driven by signs.
\index{block principal pivoting}

The second reading is the one that supports a theorem, so we develop it.

\section{Finite termination for bound constraints}

Recall the bound-constrained problem~\eqref{eq:bqp} and its free set $\Free$. The
loop solves $A_\Free x_\Free = b_\Free$, drops free variables that came back
negative, admits bound variables whose reduced gradient is negative, and repeats.

Block pivots are fast, but---like any all-at-once exchange---they can cycle: a batch
drop can overshoot the support and a later batch add restore it, revisiting a free
set already seen. The remedy is due to Júdice and Pires \cite{judice1994}: run the
batch exchange as a \emph{fast path} guarded by an anti-cycling counter, and fall
back to a single lowest-index pivot when the fast path stalls. That single pivot is
Murty's least-index principal pivoting method \cite{murty1988}, the LCP analogue of
Bland's rule, and it cannot cycle.

\begin{algorithm}[ht]
\caption{Guarded block principal pivoting}\label{alg:elim}
\begin{algorithmic}[1]
\Require Inner solver returning $x_{\Free}$ for a free set $\Free$; SPD $A$ (or its
         mat-vec) and $b$; tolerance $\varepsilon$; patience $p_{\max}$ (default $3$)
\Ensure The minimiser $x$ of~\eqref{eq:bqp}
\State $\Free \leftarrow \{1,\ldots,n\}$;\; $\bar N \leftarrow n+1$;\; $p \leftarrow p_{\max}$
\Loop
  \State solve $A_\Free x_\Free = b_\Free$;\; $x_i \leftarrow 0$ for $i \notin \Free$
  \State $s \leftarrow Ax - b$
    \hfill\textit{(reduced gradient)}
  \State $\mathcal{D} \leftarrow \{i \in \Free : x_i < -\varepsilon\}$;\;
         $\mathcal{V} \leftarrow \{i \notin \Free : s_i < -\varepsilon\}$
    \hfill\textit{(primal, dual violators)}
  \State $N \leftarrow |\mathcal{D} \cup \mathcal{V}|$
  \If{$N = 0$}
    \State \textbf{break} \hfill\textit{(KKT~\eqref{eq:lcp} holds; globally optimal)}
  \ElsIf{$N < \bar N$ \textbf{ or } $p > 0$}
    \State \textbf{if} $N < \bar N$ \textbf{ then } $\bar N \leftarrow N$;\; $p \leftarrow p_{\max}$
           \textbf{ else } $p \leftarrow p - 1$
    \State $\Free \leftarrow (\Free \setminus \mathcal{D}) \cup \mathcal{V}$
      \hfill\textit{(batch exchange)}
  \Else
    \State $i^\star \leftarrow \min\{i : i \in \mathcal{D}\cup\mathcal{V}\}$;
           toggle $i^\star$
      \hfill\textit{(Bland least-index single pivot)}
  \EndIf
\EndLoop
\end{algorithmic}
\end{algorithm}

\begin{theorem}[Finite convergence]
\label{thm:termination}
\index{finite termination!block pivoting}
Let $f(x) = \tfrac12 x\T A x - b\T x$ with $A \succ 0$. For any patience budget
$p_{\max} \in \mathbb{N}$, Algorithm~\ref{alg:elim} terminates in finitely many outer
iterations and returns the unique global minimiser of $\min_{x\ge0} f(x)$,
\emph{even under degeneracy}.
\end{theorem}

\begin{proof}
\emph{Step 1 (each inner solve is well posed).} $A \succ 0$ makes every principal
submatrix $A_\Free$ SPD (Proposition~\ref{prop:spdp}), so the inner system has a
unique solution.

\emph{Step 2 (termination implies optimality).} Free-set stationarity gives
$s_i = 0$ on $\Free$. The loop exits only when all free $x_i \ge 0$ and all bound
$s_i \ge 0$, which with complementarity is exactly~\eqref{eq:lcp}. Strict convexity
makes these sufficient (Proposition~\ref{prop:sufficient}). It therefore suffices to
show the loop cannot run forever.

\emph{Step 3 (the fallback cannot cycle).} Call a maximal block of consecutive
\textbf{else} iterates a \emph{fallback run}. Within a run $\bar N$ is constant and
each iterate performs one principal pivot on the lowest-indexed violator. This is
Murty's least-index method for the $P$-matrix LCP of Step~2, which generates a
sequence of complementary bases in which none recurs and which reaches the solution
in finitely many pivots \cite{murty1988}. A fallback run is an initial segment of
that sequence, so each run is finite wherever it starts.

\emph{Step 4 (finitely many batch steps and runs).} $\bar N$ is non-increasing and
strictly decreases exactly when the progress branch fires, so that branch fires at
most $n+1$ times. Each remaining fast-path step decrements $p$, and $p$ resets only
on progress, so at most $p_{\max}$ such steps occur between progress events.
A fallback run ends at termination or at a progress step, and there are at most
$n+1$ progress steps. Hence at most $(n+1) + (n+2)p_{\max}$ fast-path steps and at
most $n+2$ fallback runs.

\emph{Step 5.} Their sum is finite. A crude ceiling is the number of complementary
bases, $2^n$.
\end{proof}

The $2^n$ ceiling establishes \emph{finiteness} only; it is not an efficiency
estimate. The guarantee is nevertheless unconditional---the least-index fallback is
what removes the non-degeneracy hypothesis that earlier active-set arguments need,
and it is the fallback, not the batch exchange, that the proof rests on.

In practice the fallback is rarely reached. A batch step fails to reduce $N$ only
when a dropped index is immediately re-added or vice versa, which forces $s_i = 0$
exactly for some toggled $i$---a codimension-one and hence non-generic condition
\cite{judice1994}. So block pivoting converges in a handful of outer steps on
typical data with the fast path alone certifying optimality, and the guarantee is
\emph{dormant} on generic data yet genuinely necessary: on adversarial
anti-correlated designs the unguarded batch path provably cycles and only the
fallback terminates it.

\section{Why the theorem does not survive general constraints}
\label{sec:no-theorem}

Now the negative result, which is the chapter's point.

Consider $\min_{x\ge0}\tfrac12 x\T G x - a\T x$, so $C = I$ and $m = n$. The system
solved on a candidate set is
\begin{equation}
  \bigl(I_{\cdot\Active}\T G^{-1} I_{\cdot\Active}\bigr)\lambda_\Active
  = \bigl(G^{-1}\bigr)_{\Active\Active}\lambda_\Active
  = b_\Active - x_{u,\Active},
  \label{eq:principal}
\end{equation}
whose coefficient matrix is a \emph{principal submatrix} of $G^{-1}$. Since
$G^{-1} \succ 0$, every principal submatrix of it is positive definite, so $G^{-1}$
is a $P$-matrix: every principal pivot is well defined, \emph{whatever the guess}.
That is the hypothesis of Theorem~\ref{thm:termination}.

None of it is available for general $C$. The matrix $C_\Active\T G^{-1} C_\Active$
is positive definite exactly when $C_\Active$ has full column rank, and that is a
property of the guess, not of the data: a guessed set may name $k > n$ constraints,
or $k \le n$ linearly dependent ones, and then the pivot is undefined. There is no
$P$-matrix to appeal to, and correspondingly no theorem to inherit.

Contrast Proposition~\ref{prop:rank}: the dual method's step rules make dependence
\emph{unreachable}, so no rank test is needed anywhere. Guessing forfeits that
protection. The exposure is structural, not a missing anti-cycling rule, and no
tie-breaking rule repairs it, because the failure is not a cycle but an undefined
step.

\begin{remark}[Is this a real hazard or a technicality?]
\label{rem:93}
It is easy to dismiss, and the evidence says one should not. Over six hundred
random instances across four constraint families---well-posed problems of the kind
this literature routinely tests on---the rank failure \emph{never occurs}, which is
precisely why it can be overlooked. Duplicating columns of $C$, as any
programmatically assembled model may do (the same sector cap entered twice by two
model components; a constraint generated once per group and once globally), raises
it to $93\%$ of instances, and the certified fraction then collapses from $100\%$ to
nothing.

The lesson generalises past this solver: \emph{a hazard that never fires on the
family you test on is not thereby absent}. It is invisible. Test families in a
literature tend to share the structure that makes the literature's methods look
good, and finding out which structure that is, is part of reading a paper.
\end{remark}

\section{The certificate, again}

Because the method may fail---and may fail by stabilising on a set that is simply
not optimal---the trade is sound only because the answer is checked. Here
Proposition~\ref{prop:sufficient} carries the argument: a candidate $(x,\lambda)$
satisfying the tolerance test~\eqref{eq:certificate} \emph{is} the answer to that
tolerance. A candidate that fails is discarded and the exact walk of
Algorithm~\ref{alg:gi} runs instead, so guessing can never degrade the answer: it
produces the minimiser or nothing.

\begin{remark}[Two tolerances with very different jobs]
\label{rem:twotol}
The tolerance in the repair rule~\eqref{eq:repair} decides which set is tried next.
A bad choice there costs a repair or a fallback and can \emph{never} cost
correctness, so it is not delicate.

The tolerance in the certificate~\eqref{eq:certificate} is the only thing standing
between a non-optimal point and the caller, and is therefore the one to reason
about. It is nonetheless not delicate either, for the bimodality reason of
Chapter~\ref{ch:certificate}: candidates built from well-conditioned working sets
satisfy it to a few multiples of machine epsilon, and candidates built on sets that
are not have no reason to come close.

Learning to tell these two kinds of tolerance apart is worth more than any
particular tolerance value. Chapter~\ref{ch:numerics} makes it a general principle.
\end{remark}

The implementation treats the Cholesky factorisation of
$C_\Active\T G^{-1}C_\Active$ as the rank test---a dependent guess makes it fail,
rather than returning something plausible---and the least-index rule as a way of
making progress where block exchange oscillates, not as a termination proof.
Neither substitutes for the guarantee. The certificate does.

\section{Inexact inner solves}

Theorem~\ref{thm:termination} assumes each inner call returns the \emph{exact}
minimiser on the free set. Chapter~\ref{ch:krylov} will replace that call with
conjugate gradients, which returns an iterate accurate only to its residual
tolerance. Does the guarantee survive?

The question is not rhetorical: a tolerance-level solve can flip a sign test, and a
flipped sign test reintroduces exactly the cycling the fallback exists to prevent.
The answer, developed as Lemma~\ref{lem:inexact} in Chapter~\ref{ch:krylov}, is that
it does survive, provided the inner residual is tied to a quantity called the
\emph{decision margin}: the smallest nonzero magnitude any primal or dual sign test
encounters along the exact trajectory. Once the residual is small against that
margin, the inexact loop makes the identical sign decisions as the exact one and
therefore visits the identical free sets, so the termination count transfers
decision for decision.

That is a good example of the right way to compose two guarantees. One does not
prove a new termination theorem for the inexact method; one proves that the inexact
method \emph{is} the exact method, on the trajectory in question.

\begin{notes}
The primal--dual active set strategy as a semismooth Newton method is
\cite{hintermuller2002}; block principal pivoting and the least-index guard are
\cite{judice1994,murty1988,kim2011}, with \cite{cottle1992} the reference for
linear complementarity generally. Theorem~\ref{thm:termination} is the Júdice--Pires
construction restated in active-set language, following \cite{schmelzer2026nncg};
no novelty is claimed for it. The rank obstruction of
Section~\ref{sec:no-theorem} and the $600$-instance and $93\%$ measurements of
Remark~\ref{rem:93} are from the Goldfarb--Idnani note \cite{schmelzer2026quadprog}.
\end{notes}

\begin{exercises}

\begin{exercise}
Show that the starting guess ``equalities plus whatever $x_u$ violates'' returns the
correct answer in one step whenever $x_u$ is feasible.
\end{exercise}

\begin{exercise}
Verify that the repair rule~\eqref{eq:repair} is a fixed point exactly at a KKT
point, i.e.\ $\Active' = \Active$ iff \eqref{eq:pfeas}--\eqref{eq:comp} hold for the
pair produced by~\eqref{eq:pdassolve}.
\end{exercise}

\begin{exercise}
In Algorithm~\ref{alg:elim}, show that $\bar N$ can decrease at most $n+1$ times,
and construct a run in which it decreases exactly twice.
\end{exercise}

\begin{exercise}
Prove that if a batch step fails to reduce $N$ then some toggled index $i$ has
$s_i = 0$ or $x_i = 0$ exactly. Explain why this makes stalling non-generic.
\end{exercise}

\begin{exercise}
Take $G = I$ and $C$ with $c_1 = c_2 = (1,1)\T$, $b_1 = b_2 = 1$, $n = 2$. Show that
the guess $\Active = \{1,2\}$ makes $C_\Active\T G^{-1} C_\Active$ singular, and that
the dual method of Chapter~\ref{ch:walking} never forms this set.
\end{exercise}

\begin{exercise}
Design an adversarial family for which the unguarded batch path cycles. (Hint: make
the dropped index's reduced gradient turn negative as soon as it is dropped, which
anti-correlated columns achieve.)
\end{exercise}

\begin{exercise}[Implementation]
Implement Algorithm~\ref{alg:elim} with a direct free-set solve. On random SPD
problems with a planted non-negative optimum, record the number of outer steps as
$n$ grows from $50$ to $2000$. Confirm it is essentially flat. Then disable the
fallback ($p_{\max} = \infty$) and run on the adversarial family of Exercise~6 to
observe a cycle.
\end{exercise}

\begin{exercise}[Implementation]
Instrument your solver from Chapter~\ref{ch:walking} to accept a guessed working set
as a fast path, falling back to the exact walk when the certificate fails. Measure
the fraction of random instances on which the fast path succeeds, and the speedup
when it does. Then duplicate a random column of $C$ and repeat, reproducing
Remark~\ref{rem:93}.
\end{exercise}

\end{exercises}
